% Options for packages loaded elsewhere
\PassOptionsToPackage{unicode}{hyperref}
\PassOptionsToPackage{hyphens}{url}
%
\documentclass[
  10pt,
]{article}
\usepackage{amsmath,amssymb}
\usepackage{iftex}
\ifPDFTeX
  \usepackage[T1]{fontenc}
  \usepackage[utf8]{inputenc}
  \usepackage{textcomp} % provide euro and other symbols
\else % if luatex or xetex
  \usepackage{unicode-math} % this also loads fontspec
  \defaultfontfeatures{Scale=MatchLowercase}
  \defaultfontfeatures[\rmfamily]{Ligatures=TeX,Scale=1}
\fi
\usepackage{lmodern}
\ifPDFTeX\else
  % xetex/luatex font selection
\fi
% Use upquote if available, for straight quotes in verbatim environments
\IfFileExists{upquote.sty}{\usepackage{upquote}}{}
\IfFileExists{microtype.sty}{% use microtype if available
  \usepackage[]{microtype}
  \UseMicrotypeSet[protrusion]{basicmath} % disable protrusion for tt fonts
}{}
\makeatletter
\@ifundefined{KOMAClassName}{% if non-KOMA class
  \IfFileExists{parskip.sty}{%
    \usepackage{parskip}
  }{% else
    \setlength{\parindent}{0pt}
    \setlength{\parskip}{6pt plus 2pt minus 1pt}}
}{% if KOMA class
  \KOMAoptions{parskip=half}}
\makeatother
\usepackage{xcolor}
\usepackage[margin=0.85in]{geometry}
\usepackage{longtable,booktabs,array}
\usepackage{calc} % for calculating minipage widths
% Correct order of tables after \paragraph or \subparagraph
\usepackage{etoolbox}
\makeatletter
\patchcmd\longtable{\par}{\if@noskipsec\mbox{}\fi\par}{}{}
\makeatother
% Allow footnotes in longtable head/foot
\IfFileExists{footnotehyper.sty}{\usepackage{footnotehyper}}{\usepackage{footnote}}
\makesavenoteenv{longtable}
\setlength{\emergencystretch}{3em} % prevent overfull lines
\providecommand{\tightlist}{%
  \setlength{\itemsep}{0pt}\setlength{\parskip}{0pt}}
\setcounter{secnumdepth}{-\maxdimen} % remove section numbering
\ifLuaTeX
  \usepackage{selnolig}  % disable illegal ligatures
\fi
\IfFileExists{bookmark.sty}{\usepackage{bookmark}}{\usepackage{hyperref}}
\IfFileExists{xurl.sty}{\usepackage{xurl}}{} % add URL line breaks if available
\urlstyle{same}
\hypersetup{
  pdftitle={Recoverable Feature Routing under Distribution Shift: Dormant-feature adaptive re-weighting prototype},
  pdfauthor={Mykyta Siryk},
  hidelinks,
  pdfcreator={LaTeX via pandoc}}

\title{Recoverable Feature Routing under Distribution Shift:
Dormant-feature adaptive re-weighting prototype}
\author{Mykyta Siryk}
\date{4 October 2026, 18:28 UTC}

\begin{document}
\maketitle

\hypertarget{abstract}{%
\subsubsection{Abstract}\label{abstract}}

A 32-dimensional synthetic task tests whether an input-conditioned judge
can reactivate a normally dormant feature after globally strong evidence
fails. The recovered conversation reports 59.27\% linear accuracy,
52.47\% MLP accuracy, 80.10\% zero-shot adaptive accuracy, and 98.00\%
four-shot adaptive accuracy. The gate assigned to feature 23 reportedly
rises from 1.41 to 3.59. The raw result bundle was not recovered, so the
numerical record is reported-only. The prototype motivates an
image-derived falsification environment.

\hypertarget{historical-manuscript-status}{%
\subsubsection{Historical Manuscript
Status}\label{historical-manuscript-status}}

This is the cumulative scientific manuscript at release \textbf{v0.0.0}.
It includes only experiments completed by 4 October 2026, 18:28 UTC;
later findings are intentionally absent.

\hypertarget{introduction}{%
\subsubsection{1. Introduction}\label{introduction}}

Neural classifiers usually prioritize features that are most predictive
across their training distribution. That behavior is efficient when
feature reliability is stable, but brittle when a dominant cue fails
while lower-priority evidence remains informative. The motivating
deepfake-detection case is a model that trusts one family of texture,
frequency, chromatic, geometric, or localized facial evidence until
eyewear, occlusion, compression, or generator shift changes its
reliability.

The desired response is not necessarily permanent retraining. A system
should reconsider which evidence it trusts for the affected regime while
preserving its established behavior elsewhere. This manuscript records a
sequence of controlled falsification experiments designed to separate
apparent adaptation from actual access to neglected evidence.

Each tagged version is a historical scientific snapshot. Negative
results, implementation failures, and abandoned explanations are
retained because they determine why the next experiment exists.

\hypertarget{related-work}{%
\subsubsection{2. Related Work}\label{related-work}}

Mixture-of-experts systems use learned gates for input-dependent expert
selection {[}1,2{]}. MAML and fast-weight systems formalize rapid
adaptation {[}3,4{]}, while test-time adaptation and continual learning
expose complementary distribution-shift and stability problems
{[}5,6{]}. The present sequence isolates a narrower question: after a
learned evidence policy becomes unreliable, does useful backup evidence
remain both represented and reachable by a bounded adaptation process?

\hypertarget{dormant-feature-prototype}{%
\subsubsection{3. Dormant-feature
Prototype}\label{dormant-feature-prototype}}

\hypertarget{hypothesis}{%
\paragraph{3.1 Hypothesis}\label{hypothesis}}

The initial hypothesis was that an input-conditioned judge could
recognize that ordinary evidence had become unreliable and rapidly
activate a feature that was globally weak but locally decisive. This was
not posed as ordinary feature selection: the target behavior was
conditional reconsideration of evidence after an unseen transformation.

\hypertarget{construction}{%
\paragraph{3.2 Construction}\label{construction}}

The binary task used 32-dimensional samples. Features 0--5 were globally
predictive. Features 20--23 were normally noise. Three transformation
families appeared during training. A held-out transformation corrupted
strong features 0 and 5 while making feature 23 predictive. Linear and
multilayer baselines were compared with zero-shot adaptive gating and a
few-shot gate-offset update.

\hypertarget{recovered-results}{%
\paragraph{3.3 Recovered results}\label{recovered-results}}

\begin{longtable}[]{@{}lrl@{}}
\toprule\noalign{}
Method & Held-out accuracy & Evidence status \\
\midrule\noalign{}
\endhead
\bottomrule\noalign{}
\endlastfoot
Linear baseline & 0.5927 & reported-only \\
MLP baseline & 0.5247 & reported-only \\
Adaptive zero-shot & 0.8010 & reported-only \\
Adaptive four-shot & \textbf{0.9800} & reported-only \\
\end{longtable}

The gate assigned to feature 23 reportedly increased from 1.41 to 3.59
after four-shot adaptation. These values were recovered from the
conversation, not an original result archive, and therefore are not
treated as independently verified measurements.

\hypertarget{scientific-decision}{%
\paragraph{3.4 Scientific decision}\label{scientific-decision}}

The prototype showed feasibility in a deliberately favorable world. It
did not establish robustness, image-level applicability, or recovery
under a genuinely difficult failure. The next experiment therefore moved
to explicit image-derived evidence families and a held-out eyewear
intervention.

\hypertarget{interpretation-at-this-milestone}{%
\subsubsection{Interpretation at This
Milestone}\label{interpretation-at-this-milestone}}

A dormant feature can be reactivated in a favorable construction, but
the numerical evidence remains reported-only. The next stage must
introduce image-derived evidence and a meaningful held-out intervention.

\hypertarget{reason-for-the-next-stage}{%
\subsubsection{Reason for the Next
Stage}\label{reason-for-the-next-stage}}

The limitations and unresolved mechanism at this milestone directly
motivate Experiment 1; the next tagged manuscript records that test
rather than silently rewriting this version.

\hypertarget{limitations-and-threats-to-validity}{%
\subsubsection{Limitations and Threats to
Validity}\label{limitations-and-threats-to-validity}}

The controlled worlds are binary, visually regular, and explicitly
factorized. The support sets are labeled, only five principal seeds are
used after the first stages, and the synthetic glasses intervention may
be easier to recognize than natural distribution shifts.
Best-operating-point comparisons may benefit from multiple testing.
Experiment 0 lacks raw result artifacts, while Experiment 6 lacks
completed results. Executable smoke tests establish software integrity,
not reproduction of full-scale scientific outcomes.

\hypertarget{reproducibility-and-provenance}{%
\subsubsection{Reproducibility and
Provenance}\label{reproducibility-and-provenance}}

The accompanying repository separates byte-preserved historical source
and results from repaired runnable snapshots. Each runnable stage
provides a \texttt{uv} environment, exact entrypoint, bounded smoke
test, and validation record. Reconstructed proposal, implementation,
results, and analysis commits preserve why the sequence changed.
Numerical claims come from stored outputs unless explicitly marked
reported-only or results-missing.

\hypertarget{disclosure-of-ai-assistance}{%
\subsubsection{Disclosure of AI
Assistance}\label{disclosure-of-ai-assistance}}

ChatGPT collaborated in experimental ideation, implementation drafting,
debugging, interpretation, manuscript preparation, and archival
reconstruction. Git commits credit ChatGPT through co-author trailers.
The exact ChatGPT model identifiers and reasoning-effort settings used
during the original 4 October 2026 experiment conversations were not
exposed in the recovered metadata and have not been inferred. Mykyta
Siryk remains responsible for reviewing the code, evidence, claims, and
manuscript.

\hypertarget{references}{%
\subsubsection{References}\label{references}}

{[}1{]} N. Shazeer, A. Mirhoseini, K. Maziarz, A. Davis, Q. V. Le, G. E.
Hinton, and J. Dean, ``Outrageously Large Neural Networks: The
Sparsely-Gated Mixture-of-Experts Layer,'' ICLR, 2017.

{[}2{]} W. Fedus, B. Zoph, and N. Shazeer, ``Switch Transformers:
Scaling to Trillion Parameter Models with Simple and Efficient
Sparsity,'' Journal of Machine Learning Research, vol.~23, no. 120,
pp.~1--39, 2022.

{[}3{]} C. Finn, P. Abbeel, and S. Levine, ``Model-Agnostic
Meta-Learning for Fast Adaptation of Deep Networks,'' Proceedings of the
34th International Conference on Machine Learning, PMLR 70,
pp.~1126--1135, 2017.

{[}4{]} J. Ba, G. Hinton, V. Mnih, J. Z. Leibo, and C. Ionescu, ``Using
Fast Weights to Attend to the Recent Past,'' Advances in Neural
Information Processing Systems, 2016.

{[}5{]} D. Wang, E. Shelhamer, S. Liu, B. Olshausen, and T. Darrell,
``Tent: Fully Test-Time Adaptation by Entropy Minimization,''
International Conference on Learning Representations, 2021.

{[}6{]} J. Kirkpatrick et al., ``Overcoming Catastrophic Forgetting in
Neural Networks,'' Proceedings of the National Academy of Sciences,
vol.~114, no. 13, pp.~3521--3526, 2017.

\end{document}
